\section{数据划分、泄漏与分布偏移}
\label{sec:13-Machine-Learning/01-Learning-Problems-and-Evaluation/04}

\subsection{测试集 98.7\%，上线 61.3\%}

一家医学影像公司用三年间五家合作医院的数据训练肺结节检测模型，超过 120 万张胸部 CT，年龄、性别、地域都覆盖得不错。数据随机打乱后按八一一切成训练、验证、测试三份。训练曲线平滑，交叉熵从 0.52 降到 0.031，验证准确率停在 98.7\%，AUC 0.996。团队按器械监管要求锁死权重和预处理管道，封装成不可再训练的推理服务，部署到第六家医院。

第一个月的报告回来：4800 张日常片子上灵敏度只有 47\%，准确率跌到 61.3\%。更刺眼的是，对某家医院两百多张确认有结节的片子，模型只标出不到 15\%。而原来那五家医院当月新拍的片子，成绩仍稳在 96\% 以上。

数据量不是问题，收敛也不是问题——训练与验证误差的差距始终在 0.5 个百分点以内，换随机种子结果高度一致。问题出在信息流：训练脚本从数据里取走的信息比团队以为的多，而多出来的那部分在新医院并不存在。

\subsection{泛化是关于未见数据的断言，所以划分必须模拟``未见''}

测试集存在的唯一理由，是替未来的数据先站一次岗。于是切分方式必须回答一个更具体的问题：你说的``未见''，到底指未见过的哪种东西——未见过的一行记录，未见过的一段时间，还是未见过的一位患者、一家医院。同一批数据按三种方式切，模拟出来的是三种完全不同的场景。

\begin{center}
\begin{tikzpicture}[>=stealth]

\foreach \i/\p in {0/1,1/1,2/1,3/2,4/2,5/3,6/3,7/3,8/4,9/4,10/5,11/5}
  \node[draw, minimum size=0.5cm, inner sep=0, font=\scriptsize] at (\i*0.62,2.05) {\p};
\node[anchor=east, font=\scriptsize] at (-0.45,2.05) {记录（按时间排序，数字为患者编号）};

\foreach \i in {0,...,11}
  \draw[fill=black!10] (\i*0.62-0.28,1.0) rectangle (\i*0.62+0.28,1.44);
\foreach \i in {1,4,7,10}
  \draw[fill=black!55] (\i*0.62-0.28,1.0) rectangle (\i*0.62+0.28,1.44);
\node[anchor=east, font=\scriptsize] at (-0.45,1.22) {随机划分};
\node[anchor=west, font=\scriptsize] at (7.55,1.22) {患者 1 同时落在两边};

\foreach \i in {0,...,11}
  \draw[fill=black!10] (\i*0.62-0.28,0.34) rectangle (\i*0.62+0.28,0.78);
\foreach \i in {9,10,11}
  \draw[fill=black!55] (\i*0.62-0.28,0.34) rectangle (\i*0.62+0.28,0.78);
\node[anchor=east, font=\scriptsize] at (-0.45,0.56) {按时间划分};
\node[anchor=west, font=\scriptsize] at (7.55,0.56) {不许用未来预测过去};

\foreach \i in {0,...,11}
  \draw[fill=black!10] (\i*0.62-0.32,-0.32) rectangle (\i*0.62+0.32,0.12);
\foreach \i in {8,9,10,11}
  \draw[fill=black!55] (\i*0.62-0.32,-0.32) rectangle (\i*0.62+0.32,0.12);
\node[anchor=east, font=\scriptsize] at (-0.45,-0.1) {按患者分组划分};
\node[anchor=west, font=\scriptsize] at (7.55,-0.1) {整位患者留出};

\node[anchor=east, font=\scriptsize] at (-0.45,-0.85) {同一批数据，三种``未见''};
\node[anchor=east, font=\scriptsize] at (5.35,-0.85) {深色为测试集};
\draw[->] (5.58,-0.85) -- (5.58,-0.40);

\end{tikzpicture}
\end{center}

\emph{切分方式不是工程细节，它规定了测试分数在回答哪一种泛化问题。}

上图第一行就是开头那家公司的做法。它模拟的是``同一批医院里未见过的一张片子''，而部署时遇到的是``未见过的一家医院''。两个问题的答案相差 37 个百分点，中间没有任何一行代码出错。

\subsection{训练、验证和测试承担三种不同的信息角色}

三者的区别不在文件名，而在信息有没有参与过决定。训练集拟合参数；验证集参与选择——模型结构、特征集合、超参数、早停时机、决策阈值，以及最终在一堆候选流程里挑出一个；测试集只在所有选择完成之后用一次，给这套已经固定的流程一个不含选择偏差的成绩。

选择本身会制造乐观。设候选流程为 \(\mathcal A_1,\ldots,\mathcal A_m\)，用验证集挑出

\[
\widehat j=\argmin_j\widehat R_{\mathrm{val}}(\mathcal A_j),
\]

那么 \(\widehat R_{\mathrm{val}}(\mathcal A_{\widehat j})\) 已经不再无偏。设想 100 个候选的真实泛化能力完全一样，验证集上的高低纯属抽样波动，你挑走的是运气最好的那个。\(m\) 越大，挑到``看着好、实际平庸''的概率越高。独立测试集的作用，就是在 \(\widehat j\) 定下来之后重新抽一次样。

但测试集的独立性是消耗品。看过测试误差再回去改特征、换架构、调阈值，测试数据就参与了一次决定；第二次看到的数字不再是对未来的估计，而是对过去决定的反馈。这种消耗不留痕迹：文件没变，切分没变，不会有日志写``测试集独立性已于第七次查看后耗尽''，它只是悄悄退化成又一个验证集。

\subsection{切分单位必须匹配独立单位}

按行随机切分，只在每行近似独立、且部署时也是同分布随机抽样的前提下才合理。同一患者的多张影像共享解剖结构、钙化模式、摆位习惯和伪影；让这些影像分散在训练和测试两侧，模型就有机会靠``这是 3721 号患者的左肺上叶''去猜标签。测试分数很高，证明的是它学会了认人。

时间方向是另一条不可逾越的边界。随机打乱之后，完全可能用十二月的数据训练、用六月的数据测试。哪怕标签从未直接进入特征，只要预处理动用了全时间范围——比如用十二月的均值和标准差去标准化六月的数据，或者用整段时间的相关性筛特征——未来信息就已经穿越了时间边界。

群组结构同样如此：用户、设备、地点、实验批次。模型要部署到新医院，测试集就该按医院留出；要面对新用户，就该按用户留出。用 A 医院的一部分片子训练、另一部分测试，评估的是 A 医院内部的泛化；把这个数字当成跨医院迁移能力，等于假设不同医院的片子只是同一分布的不同抽样，而这个假设本身正是要检验的东西。

\subsection{预处理是训练的一部分}

标准化、缺失值填补、主成分分析、特征选择、类别重采样，在代码里常被摆在``数据准备''阶段，看着像是发生在真正的算法之前。数学上它们都在从数据里估参数。评价的对象应当是整条训练规则

\[
\mathcal A:S_{\mathrm{train}}\longmapsto\widehat f,
\]

而不是其中最后一段。凡是在划分之前动过全量数据的步骤，都已经把测试集的信息揉进了流程。

三种最隐蔽的泄漏都长在这里。特征选择泄漏是先用全部数据算特征与标签的相关性、选出最相关的一百个，再去做交叉验证——验证折的标签已经决定了哪些特征被留下，那些偶然相关在新数据上并不存在。标准化泄漏是先用全部数据算均值方差、再切分，标准化参数``知道''测试集的数值范围，表现为测试残差系统性偏小。降维泄漏是先在全部数据上做主成分分析再切分，模型在测试集上看到的是一个被测试集自己塑形过的表示空间。

这类错误的隐蔽性在于程序完全正常。对全部数据调用一次 \texttt{fit\_transform}，输出形状正确、数值合理、没有警告，没有变量被写错，没有维度不匹配，没有梯度爆炸。错误发生在信息流，不在算术表达式里。正确的做法是把 \texttt{fit} 关进训练折：每一折内部只用训练部分估计变换参数，再用同一组参数 \texttt{transform} 验证折；最终为部署在全部训练数据上重新 \texttt{fit} 也可以，但测试集从预处理到预测的整条路径上，只能被 \texttt{transform}。

\subsection{交叉验证估计的是一条训练规则，不是最终那个模型}

\(K\) 折交叉验证把数据分成 \(K\) 份，轮流用 \(K-1\) 份训练、剩下一份评价，再取平均。它比单次留出更省数据，适合比较流程和挑超参数。

但它估的不是``用全部 \(n\) 个样本训出来的那个模型''的表现，而是``用大约 \((K-1)n/K\) 个样本训练的这条规则''在当前切分机制下的平均表现。样本量不大、学习曲线还在爬升时，这个差别会带来系统性低估：交叉验证的平均误差略高于最终全数据模型的真实误差，样本越少、曲线越陡，差得越多。这一点在\hyperref[sec:11-Mathematical-Statistics/06-Resampling-and-Model-Selection/03]{``交叉验证评估的是完整训练流程''}里有更细的展开，那里也说明了为什么被评估的对象必须是整条流程而不是某一个拟合结果。

若超参数也靠交叉验证来选，而又想给整个选择过程一个诚实的成绩，就需要嵌套交叉验证：外层把数据切成训练折与测试折，内层只在外层训练折上再做一次交叉验证来挑超参，外层测试折在此期间完全不被触碰，包括预处理、特征选择、超参搜索和早停判断。时间序列的交叉验证要保持方向，滚动窗口或扩展窗口都行，前提是训练段永远早于评价段。群组数据则按群组整块留出，同一位患者的片子不允许分散到不同折。折数只是计算预算，按什么单位切，才决定了这次交叉验证在模拟哪一种泛化。

\subsection{三种分布偏移改变的是分解中的不同因子}

联合分布有两种分解方式：

\[
P(X,Y)=P(Y\given X)\,P(X)=P(X\given Y)\,P(Y).
\]

协变量偏移改的是 \(P(X)\)，而暂时假设 \(P(Y\given X)\) 不变：新医院的患者年龄构成、设备型号、扫描参数都不同，但给定完全相同的影像表现，结节的判断标准与旧医院一致。标签偏移改的是 \(P(Y)\)，暂时假设 \(P(X\given Y)\) 不变：疫情期间肺炎比例骤升，而``肺炎''这一类的影像表现没变，只是它在人群里出现得更频繁。概念漂移改的是 \(P(Y\given X)\) 本身：同样的影像在新的诊断标准或阅片习惯下对应不同结论，输入到输出的映射变了，这是最难修的一种。

要注意这三个名字说的都是``哪个因子被假设不变''，不是能从两张直方图上读出来的事实。\(P(X)\) 与 \(P(Y)\) 同时变化时，光看数据分不清是协变量偏移、标签偏移还是二者叠加，得靠额外的结构假设或领域知识。

协变量偏移在训练分布覆盖了部署分布的支撑集、且密度比能可靠估计时，可以用重要性加权修正：

\[
\widehat R_{\mathrm{weighted}}(f)=\frac1n\sum_{i=1}^n\frac{p_{\mathrm{deploy}}(x_i)}{p_{\mathrm{train}}(x_i)}\,\ell\bigl(f(x_i),y_i\bigr).
\]

部署时常见而训练时罕见的区域权重大于 1，损失被放大；反之被降权。两个前提都很硬。其一，训练分布在部署分布支撑集的每一点上密度为正；新医院用的设备型号在训练数据里从未出现，对应特征空间中一整块新区域，分母为零，权重无定义。其二，高维密度比的估计本身就是难题，估计误差可能比原问题还大。

还有一类偏移是模型自己造成的。推荐系统决定用户看到什么，被推荐内容的点击率不等于用户自主浏览时的点击率；风控模型拦下的交易不再产生标签，那些模式在后续数据里永远停在``已拦截''，无从确认是欺诈还是正常；强化学习的状态分布是策略的函数，不是环境的固定属性。这些场景里数据不再是被动环境的静态抽样，评价必须把行动与反馈的回路算进去。

\subsection{测试分数本身也是随机量}

即使测试集完全独立、分布也对得上，报告出来的经验风险仍带着抽样误差。逐样本损失方差有限且样本近似独立时，均值的标准误按 \(1/\sqrt n\) 收缩；错误率接近 \(p\) 时，独立 Bernoulli 近似给出标准误 \(\sqrt{p(1-p)/n}\)。取 \(n=1000\)、\(p=0.05\)，标准误约 0.7 个百分点，于是报告 5.0\% 错误率时，真实值落在 3.6\% 到 6.4\% 之间都不奇怪。测试集扩到 10000，标准误降到约 0.22 个百分点——十倍样本换来三倍精度，小数点后第三位在多数样本量下不携带任何可靠信息。

这个公式有三处会失效。样本相关时有效样本量小于观测数，标准误被低估，同一位患者的多次测量就是典型；损失长尾时方差估计被少数极端样本主导，有限方差这个前提本身就悬；同时比较多个模型并挑出最好的那个，挑选动作又把乐观性带了回来，即使每个测试误差单独无偏，最小值也系统性偏低。

比较两个模型时，更有效的做法是在同一组测试样本上算逐样本损失差 \(\ell(f_1(x_i),y_i)-\ell(f_2(x_i),y_i)\)，再对这列差值求均值与标准误。难例对两个模型造成的共同波动在差分时抵消，差的方差小于两个方差之和，同样的样本量能分辨更小的差距。

\subsection{一条可以照着做的评估流程}

把上面这些拼起来，一次可信的评估大致按这个顺序走。先确定``未见''的单位：部署时遇到的新东西是新时间、新实体还是新场地，切分就按哪一个来；不确定时，按最保守的那个切，同时保留一份随机切分的结果作对照，两者的差距本身就是一份关于泄漏和偏移的诊断。切分之后，任何需要估参数的步骤都关在训练侧，测试侧只允许被变换。

接着确定指标，而不是接受默认值。指标要由代价和容量决定：代价不对称就别看准确率，只能复核前 200 笔就报排序顶部的精确率，需要概率去做阈值决策就加上校准误差。总体一行之外必须有分组几行，否则读者无法排除总体与分组方向相反的情形。同时定下基线——常数基线、持续性基线、单特征基线、当前线上规则——并让基线走完全相同的划分与预处理。

然后才是比较。任何差距都要配区间：对配对损失差算标准误，或者用自助法重采样测试集得到区间；区间盖过零，就当作没有结论。这一步还要记账：为了得到这个结果，测试集被看过几次，候选流程有多少个。选择次数是估计乐观性的一部分，不写下来就无法评估。

最后一步是判断离线结论到此为止还是需要在线实验。离线评估在数据被行动改变时必然失真——模型改变哪些交易被拦截、哪些内容被曝光，历史数据里就不存在反事实的标签。业务后果无法从离线标签重建（用户留存、投诉率、审核工时）、系统因素参与决定结果（延迟、界面、上下游策略）、或者模型将改变自己未来的训练数据，这三种情形下都该做在线随机实验。实验的随机化单位要与干扰单位一致：同一位用户或同一家门店内部的样本会互相影响，随机化就得抬到用户或门店一级；实验前先按预期效应量算功效，定好时长和样本量，别看到显著就停。

\subsection{评价纪律能防住什么，防不住什么}

严格划分训练、验证与测试，防的是同一份随机性被反复利用，它保证报告的数字不因选择乐观而虚高。但一个毫无泄漏的实验，仍然可能精确地回答一个错误的问题。标签定义可能与业务目标不一致——标注意义上的``结节''与临床需要干预的``可疑结节''是两个集合；训练样本可能不覆盖目标人群——基层医院的影像特征与三甲医院不同；数据生成机制可能在部署期改变——诊断指南更新之后，同样的影像学表现对应不同结论。

开头那家公司正落在这一格里。98.7\% 与 61.3\% 之间的差距不来自泄漏，切分本身没错；差距来自测试数据与训练数据出自同一批医院，而部署发生在新医院。评价设计成功地阻止了同一份随机性被用两次，却没有阻止把旧世界的规律当成新世界的保证。

所以可信评价始终靠两层检查。统计层保留独立信息、量化抽样误差、防止把偶然当规律，这是划分与交叉验证解决的事。建模层要求评价数据真的模拟部署时的对象、时间与行动，这要求你知道数据是怎么收集的、模型将被用在哪里、什么样的变化会让过去的规律在明天失效。前者让估计诚实，后者让诚实的估计仍然与即将面对的困难有关。这一单元到此立好了评价的规矩；从下一单元起，所有模型都要在这套规矩下接受比较。
